% !TEX TS-program = pdflatex
\documentclass[12pt,a4paper]{article}
\usepackage[utf8]{inputenc}
\usepackage[T1]{fontenc}
\usepackage{lmodern}
\usepackage[a4paper,margin=2.05cm]{geometry}
\usepackage{amsmath,amssymb,mathtools}
\usepackage{graphicx}
\usepackage{booktabs}
\usepackage{longtable}
\usepackage{array}
\usepackage{hyperref}
\usepackage{caption}
\usepackage{float}
\usepackage{enumitem}
\usepackage{xcolor}
\usepackage{microtype}
\hypersetup{colorlinks=true,linkcolor=blue,urlcolor=blue,citecolor=blue}
\setlength{\parskip}{0.55em}
\setlength{\parindent}{0pt}
\renewcommand{\arraystretch}{1.15}
\newcommand{\TMI}{\mathrm{TMI}}
\newcommand{\std}{\mathrm{std}}
\newcommand{\base}{\mathrm{base}}
\newcommand{\Prob}{\operatorname{Pr}}
\newcommand{\Auto}{\operatorname{Auto}}
\newcommand{\Safe}{\operatorname{Safe}}

\title{\textbf{Computational Validation of Adaptive Interface Control in TMI Quantum Systems}\\
\large Extended Technical Report on Computational Tests 1--11}
\author{Salkutsan Aleksey}
\date{May 6, 2026}

\begin{document}
\maketitle

\begin{abstract}
This document consolidates the computational test line for quantum-system predictions in the Theory of Manifold Interfaces (TMI). This version extends the initial Tests 1--6 and adds Tests 7--11: full density-matrix Monte Carlo validation, cross-model validation, failure-boundary mapping, adversarial failure testing, and safety-aware recovery testing.

The tested engineering core has the form
\[
I_Q^{coh-auto}\Rightarrow Loss_{coh}\downarrow\Rightarrow T_2^{eff}\uparrow\Rightarrow P_{success}\uparrow.
\]
The report does not claim a laboratory proof of the physical existence of the field $\Xi_I^{auto}$. It fixes a more modest result: in computational models, adaptive interface control can increase effective coherence and circuit success probability, but it has a domain of applicability and requires a safety-aware layer under adversarial conditions.
\end{abstract}

\tableofcontents
\newpage

\section{Status and Scientific Caution}

This document is a \textbf{computational validation report}, not a laboratory confirmation. We do not attempt to verify the entire Theory of Manifold Interfaces at once. We test one narrow consequence:
\[
\boxed{
I_Q^{coh-auto}
\Rightarrow
\gamma_{eff}\downarrow
\Rightarrow
T_2^{eff}\uparrow
\Rightarrow
P_{success}\uparrow
}
\]
where $I_Q^{coh-auto}$ is a coherence-preserving automatic interface, $\gamma_{eff}=1/T_2^{eff}$ is the effective decoherence rate, and $P_{success}$ is the probability of obtaining the correct result in a small quantum circuit.

Starting with Test 10 and Test 11, the report also fixes a negative and safety-related consequence:
\[
\boxed{\Auto(I)\not\Rightarrow \Safe(I)}
\]
and
\[
\boxed{\Auto(I)+\text{SafetyLayer}\Rightarrow \text{SaferAdaptiveControl}.}
\]
In other words, adaptivity by itself is not a guarantee of safety: under misleading feedback and high control cost, an automatic interface can lose to a fixed standard control mode.

\section{Connection to TMI Core 0.2 and TMI-Physics 2.2}

The computational test line does not replace Core 0.2 or TMI-Physics 2.2. Its correct role is that of an applied computational module:
\[
\boxed{
\text{TMI Core 0.2}
\Rightarrow
\text{TMI-Physics 2.2}
\Rightarrow
\text{Computational Validation}
}
\]

In Core 0.2, automaticity is understood as a definitional property of a valid interface:
\[
\operatorname{Valid}(I)\Rightarrow \Auto(I),
\qquad
\Auto(I)\not\Rightarrow \operatorname{AlwaysAct}(I).
\]
In TMI-Physics 2.2 this logic is projected into the physical layer:
\[
\Xi_I\in\mathfrak I_{phys}
\Rightarrow
\Auto(\Xi_I)
\Rightarrow
\Xi_I^{auto}.
\]
The computational tests examine an engineering projection of this principle: can an automatic coherence-control interface improve observable metrics in quantum-system models?

\section{Summary of Tests 1--11}

For Tests 1--4, the final column typically lists the values in the order base / standard / TMI.

\begin{longtable}{p{2.0cm}p{4.0cm}p{4.2cm}p{5.0cm}}
\toprule
Test & Meaning & Criterion & Result \\
\midrule
\endfirsthead
\toprule
Test & Meaning & Criterion & Result \\
\midrule
\endhead

Test 1 & Ramsey toy model & $T_2^{TMI}>T_2^{std}>T_2^{base}$ & 50.0 / 120.0 / 178.6 $\mu s$ \\
Test 2 & Noisy Ramsey learning & $T_2^{TMI}>T_2^{std}>T_2^{base}$ & 51.3 / 125.8 / 162.4 $\mu s$ \\
Test 3 & Small circuit & $P_{TMI}>P_{std}>P_{base}$ & 0.182 / 0.386 / 0.431 \\
Test 4 & Stabilized interface & $P_{stab}>P_{std}>P_{base}$ & 0.180 / 0.328 / 0.465 \\
Test 5 & Surrogate Monte Carlo & $E[P_{TMI}]>E[P_{std}]$ & mean $\Delta=0.0512$, $Pr=0.864$ \\
Test 6 & Ablation study & full TMI $>$ standard & full $\Delta\approx 0.0677$ \\
Test 7 & Full density-matrix MC & $E[P_{TMI}]>E[P_{std}]$ & mean $\Delta=0.0935$, $Pr=0.933$ \\
Test 8 & Cross-model validation & advantage across circuit families & mean $\Delta=0.0662$, $Pr=0.846$ \\
Test 9 & Failure-boundary map & locate weak/failure regions & failure cells $=0$, weak cells $=6$ \\
Test 10 & Adversarial failure & force $P_{TMI}\le P_{std}$ & failure cells $=58$, weak cells $=22$ \\
Test 11 & Safety recovery & reduce failure rate & $\Delta_{safety}=0.0139$, recovered cells $=31/58$ \\

\bottomrule
\end{longtable}

\section{Tests 1--6: Basic Computational Validation Line}

\subsection{Test 1: Ramsey toy model}

The coherence model is
\[
C(t)=C_0 e^{-t/T_2}.
\]
The result is
\[
T_2^{base}=50.0\,\mu s,
\quad
T_2^{std}=120.0\,\mu s,
\quad
T_2^{TMI}=178.6\,\mu s.
\]

\begin{figure}[H]
\centering
\includegraphics[width=0.86\linewidth]{figures/fig_test1_ramsey.png}
\caption{Test 1: simple Ramsey coherence model.}
\end{figure}

\subsection{Test 2: Noisy Ramsey learning}

The automatic interface must learn the protection parameter $u$ from noisy Ramsey observations. The result is
\[
T_2^{base}=51.3\,\mu s,
\quad
T_2^{std}=125.8\,\mu s,
\quad
T_2^{TMI}=162.4\,\mu s.
\]

\begin{figure}[H]
\centering
\includegraphics[width=0.86\linewidth]{figures/fig_test2_noisy_ramsey.png}
\caption{Test 2: Ramsey curves under noisy observations.}
\end{figure}

\begin{figure}[H]
\centering
\includegraphics[width=0.86\linewidth]{figures/fig_test2_learning.png}
\caption{Test 2: learning of the protection parameter $u$.}
\end{figure}

\subsection{Test 3: Small quantum circuit}

The circuit is constructed as
\[
|000\rangle\to \text{entangling circuit}\to \text{inverse circuit}\to |000\rangle.
\]
The ideal success probability is
\[
P_{success}=\Pr(000).
\]
The result is
\[
P_{base}=0.182,\quad P_{std}=0.386,\quad P_{TMI}=0.431.
\]

\begin{figure}[H]
\centering
\includegraphics[width=0.82\linewidth]{figures/fig_test3_success.png}
\caption{Test 3: success probability of a small quantum circuit.}
\end{figure}

\subsection{Test 4: Stabilized adaptive interface}

Test 4 adds memory, smoothing, step limitation, and repeated observations:
\[
\text{AutoInterface}+\text{Memory}+\text{Smoothing}+\text{StepLimit}\Rightarrow\text{StableControl}.
\]
The result is
\[
P_{base}=0.180,\quad P_{std}=0.328,
\quad P_{raw}=0.465,\quad P_{stab}=0.465.
\]

\begin{figure}[H]
\centering
\includegraphics[width=0.82\linewidth]{figures/fig_test4_success.png}
\caption{Test 4: stabilized adaptive interface improves $P_{success}$.}
\end{figure}

\begin{figure}[H]
\centering
\includegraphics[width=0.86\linewidth]{figures/fig_test4_learning.png}
\caption{Test 4: raw adaptive control jumps, while stabilized adaptive control keeps the parameter closer to the optimal regime.}
\end{figure}

\subsection{Test 5: Monte Carlo surrogate validation}

Test 5 checks statistical advantage across 500 random environments using a fast circuit-level surrogate:
\[
\mathbb E[P_{TMI}]>\mathbb E[P_{std}],
\qquad
\Pr(P_{TMI}>P_{std})>0.5.
\]
The result is
\[
\mathbb E[P_{base}]=0.245,\quad
\mathbb E[P_{std}]=0.417,\quad
\mathbb E[P_{TMI}]=0.468.
\]
\[
\mathbb E[P_{TMI}-P_{std}]=0.0512,
\quad
\Pr(P_{TMI}>P_{std})=0.864.
\]

\begin{figure}[H]
\centering
\includegraphics[width=0.82\linewidth]{figures/fig_test5_mean.png}
\caption{Test 5: mean $P_{success}$ in Monte Carlo surrogate validation.}
\end{figure}

\begin{figure}[H]
\centering
\includegraphics[width=0.82\linewidth]{figures/fig_test5_delta.png}
\caption{Test 5: distribution of the advantage $P_{TMI}-P_{std}$.}
\end{figure}

\subsection{Test 6: Ablation study}

Test 6 shows that raw adaptivity alone does not guarantee an improvement:
\[
\text{RawAdaptivity}\not\Rightarrow\text{Improvement}.
\]
The advantage appears as a property of a structured automatic interface:
\[
\text{StructuredAutoInterface}\Rightarrow\text{StableParameterTracking}\Rightarrow P_{success}\uparrow.
\]

\begin{figure}[H]
\centering
\includegraphics[width=0.90\linewidth]{figures/fig_test6_ablation.png}
\caption{Test 6: ablation study across interface-control configurations.}
\end{figure}

\begin{figure}[H]
\centering
\includegraphics[width=0.90\linewidth]{figures/fig_test6_delta_full.png}
\caption{Test 6: losses relative to the full TMI mechanism.}
\end{figure}

\section{Tests 7--8: Robustness and Cross-Model Validation}

\subsection{Test 7: Full density-matrix Monte Carlo}

Test 7 repeats the statistical validation of Test 5, but now on the full 3-qubit density-matrix toy model. This removes the main limitation of Test 5: dependence on a fast surrogate.

\[
\mathbb E[P_{base}]=0.182,\quad
\mathbb E[P_{std}]=0.322,\quad
\mathbb E[P_{TMI}]=0.415.
\]
\[
\mathbb E[P_{TMI}-P_{std}]=0.0935,
\quad
\Pr(P_{TMI}>P_{std})=0.933.
\]

\begin{figure}[H]
\centering
\includegraphics[width=0.82\linewidth]{figures/fig_test7_mean.png}
\caption{Test 7: full density-matrix Monte Carlo, mean $P_{success}$.}
\end{figure}

\begin{figure}[H]
\centering
\includegraphics[width=0.70\linewidth]{figures/fig_test7_paired.png}
\caption{Test 7: paired comparison $P_{std}$ vs $P_{TMI}$.}
\end{figure}

\subsection{Test 8: Cross-model validation}

Test 8 checks not one circuit, but four circuit families. The overall result is
\[
\mathbb E[P_{base}]=0.292,\quad
\mathbb E[P_{std}]=0.379,\quad
\mathbb E[P_{TMI}]=0.445.
\]
\[
\mathbb E[P_{TMI}-P_{std}]=0.0662,
\quad
\Pr(P_{TMI}>P_{std})=0.846.
\]

\begin{longtable}{p{5.3cm}rrrrr}
\toprule
Circuit & Base & Standard & TMI & $\Delta$ & $Pr$ \\
\midrule
\endfirsthead
\toprule
Circuit & Base & Standard & TMI & $\Delta$ & $Pr$ \\
\midrule
\endhead
balanced\_entangle\_inverse & 0.169 & 0.277 & 0.376 & 0.0994 & 1.000 \\
deep\_chain\_inverse & 0.253 & 0.322 & 0.377 & 0.0553 & 0.923 \\
ghz\_long\_idle\_inverse & 0.495 & 0.512 & 0.527 & 0.0148 & 0.462 \\
local\_mix\_inverse & 0.249 & 0.404 & 0.500 & 0.0955 & 1.000 \\
\bottomrule
\end{longtable}

\begin{figure}[H]
\centering
\includegraphics[width=0.95\linewidth]{figures/fig_test8_mean_by_circuit.png}
\caption{Test 8: mean $P_{success}$ by circuit family.}
\end{figure}

\begin{figure}[H]
\centering
\includegraphics[width=0.90\linewidth]{figures/fig_test8_delta_by_circuit.png}
\caption{Test 8: TMI advantage over standard control by circuit family.}
\end{figure}

\paragraph{Important nuance.}
For the \texttt{ghz\_long\_idle\_inverse} circuit, the mean advantage is positive but weaker and less stable than for the other circuit families. This shows that the mechanism is not equally strong for all circuit types.

\section{Tests 9--11: Applicability Boundaries and Safety Recovery}

\subsection{Test 9: Failure-boundary test}

Test 9 searches not only for success, but also for weakening zones. We vary the drift speed of the hidden optimum and the measurement noise. The failure criterion is
\[
\mathbb E[P_{TMI}-P_{std}]\le 0
\quad\text{or}\quad
\Pr(P_{TMI}>P_{std})\le 0.5.
\]

The overall result is
\[
\mathbb E[P_{TMI}-P_{std}]=0.0606,
\quad
\Pr(P_{TMI}>P_{std})=0.849.
\]
Strict failure cells found:
\[
0,
\]
weak zones found:
\[
6.
\]

\begin{figure}[H]
\centering
\includegraphics[width=0.82\linewidth]{figures/fig_test9_delta_heatmap.png}
\caption{Test 9: heatmap of the advantage $P_{TMI}-P_{std}$ over drift/noise.}
\end{figure}

\begin{figure}[H]
\centering
\includegraphics[width=0.82\linewidth]{figures/fig_test9_prob_heatmap.png}
\caption{Test 9: probability $\Pr(P_{TMI}>P_{std})$ over drift/noise.}
\end{figure}

\subsection{Test 10: Adversarial failure test}

Test 10 intentionally creates adversarial conditions:
\[
\text{fast drift}+\text{high noise}+\text{misleading feedback}+\text{high control penalty}.
\]
The goal is to obtain real cases of
\[
P_{TMI}\le P_{std}.
\]

The result is
\[
\mathbb E[P_{TMI}-P_{std}]=0.0025,
\quad
\Pr(P_{TMI}>P_{std})=0.530.
\]
Failure cells:
\[
58,
\]
weak cells:
\[
22.
\]

The key conclusion is
\[
\boxed{\text{AdaptiveInterface}\not\Rightarrow\text{SafeInterface}.}
\]

\begin{figure}[H]
\centering
\includegraphics[width=0.84\linewidth]{figures/fig_test10_failure_dist.png}
\caption{Test 10: adversarial failure distribution.}
\end{figure}

\begin{figure}[H]
\centering
\includegraphics[width=0.70\linewidth]{figures/fig_test10_failure_heatmap.png}
\caption{Test 10: number of failure cells by misleading bias and control penalty.}
\end{figure}

\subsection{Test 11: Safety-aware recovery}

Test 11 takes only the failure cells from Test 10 and checks whether the advantage can be recovered by a safety layer:
\[
\text{AutoInterface}+\text{SafetyPenalty}+\text{ConservativeFallback}+\text{ValidationChannel}\Rightarrow\text{SaferControl}.
\]

\begin{longtable}{p{3.8cm}rrrrr}
\toprule
Mode & Mean $P$ & $\Delta$ vs std & $Pr(>std)$ & Failure rate & Mean $u$ \\
\midrule
\endfirsthead
\toprule
Mode & Mean $P$ & $\Delta$ vs std & $Pr(>std)$ & Failure rate & Mean $u$ \\
\midrule
\endhead
standard & 0.2799 & 0.00000 & 0.000 & -- & 0.450 \\
adversarial\_TMI & 0.2799 & -0.00002 & 0.478 & 0.522 & 0.472 \\
safety\_v1 & 0.2939 & 0.01392 & 0.617 & 0.383 & 0.395 \\
recovery\_v2 & 0.2932 & 0.01328 & 0.579 & 0.421 & 0.386 \\
\bottomrule
\end{longtable}

The main result is
\[
\Delta_{adv}=-0.00002,
\quad
\Delta_{safety}=0.01392,
\quad
\Delta_{recovery}=0.01328.
\]
Recovery v2 recovered
\[
31/58
\]
failure cells.

\begin{figure}[H]
\centering
\includegraphics[width=0.82\linewidth]{figures/fig_test11_boxplot.png}
\caption{Test 11: safety recovery on failure zones from Test 10.}
\end{figure}

\begin{figure}[H]
\centering
\includegraphics[width=0.82\linewidth]{figures/fig_test11_failure_rate.png}
\caption{Test 11: failure-rate reduction under safety-aware control.}
\end{figure}

\section{Final Theoretical Consequences}

\subsection{Positive consequence}

Under ordinary and moderately noisy conditions, the computational line supports the engineering hypothesis
\[
\boxed{
I_Q^{coh-auto}
\Rightarrow
T_2^{eff}\uparrow
\Rightarrow
P_{success}\uparrow
}
\]

\subsection{Negative consequence}

Test 10 introduces an important boundary:
\[
\boxed{
\text{MisleadingFeedback}\land\text{HighControlPenalty}
\Rightarrow
P_{TMI}\le P_{std}
}
\]
Thus, an automatic interface can make errors if its feedback channel systematically misleads it.

\subsection{Safety-aware consequence}

Test 11 shows that a safety layer reduces the failure region, but does not eliminate it completely:
\[
\boxed{
\text{AutoInterface}+\text{SafetyLayer}
\Rightarrow
\text{SaferAdaptiveControl}
}
\]
but not
\[
\text{AutoInterface}+\text{SafetyLayer}
\Rightarrow
\text{PerfectControl}.
\]

\section{Next Steps}

The next steps are:
\begin{enumerate}[label=\arabic*.]
\item Prepare an independent implementation of Tests 7--8 in a Qiskit/Cirq-style framework.
\item Repeat the key tests on a standard noise simulator.
\item Separate the theoretical TMI-Physics 2.2 module from the computational supplement.
\item Prepare a reproducible repository structure: code, data, figures, and reports.
\item After independent simulator validation, move to a cloud-QPU protocol.
\end{enumerate}

\section{Conclusion}

Tests 1--8 show that adaptive interface control can improve $T_2^{eff}$ and $P_{success}$ in numerical models. Tests 9--11 make the line scientifically more honest: they show not only the region of success, but also the region of weakness, failure, and partial recovery through a safety-aware layer.

The final formula of the report is
\[
\boxed{
\text{Adaptive Interface Control}
\Rightarrow
\text{Measurable Computational Advantage}
}
\]
provided that the environment is observable enough, the feedback channel is stable enough, and control is constrained by safety-aware limits.

\appendix
\section{Package Contents}

The report archive contains:
\begin{itemize}
\item the PDF and LaTeX source of this report;
\item figures/ - key plots from Tests 1--11;
\item data/ - CSV results and summary files;
\item test\_archives/ - archives of individual Tests 1--11.
\end{itemize}

\end{document}
